\documentclass[runningheads]{llncs}

\usepackage[T1]{fontenc}
\usepackage[utf8]{inputenc}
\usepackage{graphicx}
\usepackage{booktabs}
\usepackage{multirow}
\usepackage{amsmath}
\usepackage{amssymb}
\usepackage{url}

\title{Rumour Detection on Twitter Cascades:\\ A Graph Neural Network Approach with Topological and XAI Analysis}
\titlerunning{Rumour Detection on Twitter Cascades}
\author{Nuno Batista \and Miguel Cabral Pinto}
\authorrunning{N. Batista and M. Cabral Pinto}
\institute{University of Coimbra, Coimbra, Portugal}

\begin{document}
\maketitle

\begin{abstract}
Automated rumour detection on social media requires classifiers that generalise
across unseen news events, yet most prior work evaluates within a single
train/test split. We study this challenge on the PHEME dataset of 6{,}178
Twitter cascade graphs spanning seven real-world events, using a strict
two-phase leave-one-event-out protocol that keeps two events entirely withheld
during hyperparameter search. We benchmark Random Forest, MLP, GCN, and GAT
under three feature ablations (full, text-only, structural-only), and complement
the classification study with a network science characterisation of cascade
topology and GNNExplainer and GAT attention interpretability analyses.
Structural-only models fall below the random lower bound on the held-out test
set (macro F1 $< 0.375$), confirming that pure topology does not generalise
across events, despite consistent structural differences between classes
(rank-biserial $r < 0.15$ throughout). Text-based models all exceed the lower
bound; GAT with text features achieves the best combined macro F1 of 0.610,
significantly outperforming flat mean-pooled MLP (McNemar $p < 0.001$).
Adding structural features to text-based models consistently degrades held-out
performance, indicating topology is event-specific rather than event-general.
GNNExplainer reveals that authority signals (follower count, PageRank) are
more important for rumour predictions than non-rumour ones, while GAT attention
concentrates on fewer nodes for rumours, consistent with an anomalous-node
detection pattern.
\end{abstract}

\section{Introduction}

Misinformation on social media poses a challenge that scales with platform size.
Vosoughi et al.~\cite{vosoughi2018spread} found that false news propagates roughly
six times faster than true news on Twitter, reaching far larger audiences before
corrections can circulate. The consequences, from electoral interference to
vaccine hesitancy, have sustained a substantial body of automated detection
research.

A Twitter reply thread is a natural directed graph: the source tweet is the root
and each reply connects to its parent, forming a cascade tree that encodes both
content and propagation dynamics. Most early detection systems treat rumour
classification as a flat text problem, aggregating per-tweet features into a
single vector and discarding the relational structure. Graph Neural Networks
(GNNs)~\cite{kipf2017semi,velickovic2018gat} make it possible to learn directly
over these cascade graphs. Prior work, notably BiGCN~\cite{bian2020bigcn} (which propagates information
bidirectionally through the cascade tree) and geometric deep learning on
propagation graphs~\cite{monti2019fakenews}, has shown that exploiting graph
structure improves over flat baselines. However, these studies
typically evaluate within a single train/test split, leaving open whether GNNs
learn event-transferable cascade signatures or whether apparent gains reflect
overfitting to training-event topology distributions.

This paper addresses three questions using the PHEME dataset~\cite{zubiaga2016pheme},
a collection of 6{,}425 Twitter cascade graphs across nine real-world news events
labelled as Rumour or Non-Rumour. First, does graph topology alone generalise
across unseen events? Second, does message passing add value beyond
mean-pooled text embeddings when node features are held constant? Third, which
features drive GNN predictions, and do these patterns differ between classes?
We adopt a strict leave-one-event-out (LOEO) evaluation protocol, keeping two
events entirely held out from hyperparameter search.

Our contributions are threefold. We characterise cascade topology across six
analyses spanning cascade size and depth, structural virality~\cite{goel2016structural},
network robustness, echo chambers, and novel branching metrics, finding consistent
but small structural differences: rumour cascades are shallower and more
broadcast-like, with lower Strahler branching complexity, yet these effects are
insufficient to drive cross-event classification (structural-only models fall below
the random lower bound). We then benchmark MLP, GCN, and GAT under three feature ablations (full,
text-only, and structural-only) with bootstrap confidence intervals and
McNemar's significance tests; GAT with text features achieves a combined macro
F1 of 0.610, significantly outperforming flat mean-pooled MLP ($p < 0.001$),
confirming that message passing provides genuine signal beyond text alone. Finally, we apply GNNExplainer~\cite{ying2019gnnexplainer} and GAT attention
analysis: GNNExplainer reveals the GCN assigns significantly higher feature
importance to authority signals (follower count, PageRank) for rumour graphs, while
GAT attention is highest at depth-2 in both classes and shows lower per-graph
entropy for rumours, indicating more concentrated attention.

\section{Dataset and Problem Formulation}

\subsection{The PHEME Dataset}

We formulate rumour detection as binary graph-level classification. Given a
cascade $\mathcal{G} = (\mathcal{V}, \mathcal{E}, \mathbf{X})$ (nodes = tweets,
edges = replies, $\mathbf{X} \in \mathbb{R}^{|\mathcal{V}| \times d}$ the node
feature matrix), the goal is to predict $y \in \{0, 1\}$ (Non-Rumour / Rumour)
for the graph as a whole.

We use the PHEME dataset~\cite{zubiaga2016pheme}, comprising nine real-world
Twitter news events with thread-level Rumour/Non-Rumour annotations. Two events
are excluded due to near-degenerate class distributions: \textit{ebola-essien}
(100\% Rumour, 14 cascades) and \textit{prince-toronto} (98.3\% Rumour, 233
cascades). A further event, \textit{ferguson}, is excluded from classifier
evaluation owing to distributional shift that substantially inflates
cross-event variance. The remaining seven events yield 6,178 graphs in total;
event-level statistics are shown in Table~\ref{tab:dataset}. The positive rate
varies from 22.0\% (charliehebdo) to 52.9\% (putinmissing), reflecting genuine
differences in how each news event unfolded on Twitter. We apply a
minimum-size filter of $|\mathcal{V}| \geq 3$, removing cascades that consist
only of a source tweet and at most one reply.

Class imbalance is handled at training time by per-event majority-class
undersampling, ensuring that no single event's distribution dominates the
training set.

As an additional contribution, we convert the PHEME annotations and graph
structures into a ready-to-use PyTorch Geometric dataset and release it
publicly on Zenodo~\cite{batista2026pheme_pyg} and Hugging
Face\footnote{\url{https://huggingface.co/datasets/NunoBatista/PHEME-Misinformation-Graphs}}
for convenient programmatic access, removing the need to parse raw JSON files
to reproduce our experiments.

\begin{table}[t]
\caption{PHEME event statistics after excluding \textit{ebola-essien} and
\textit{prince-toronto}. Events marked ($\dagger$) form the held-out test set;
($\ddagger$) is excluded from classifier evaluation (distribution shift) and
used only in the network-science analysis.}
\label{tab:dataset}
\centering
\begin{tabular}{lrrrr}
\toprule
Event & Graphs & Rumour & Non-Rumour & R\% \\
\midrule
charliehebdo$^{\dagger}$  & 2{,}079 & 458 & 1{,}621 & 22.0 \\
ferguson$^{\ddagger}$     & 1{,}143 & 284 &   859   & 24.8 \\
germanwings               &   469   & 238 &   231   & 50.7 \\
gurlitt                   &   138   &  61 &    77   & 44.2 \\
ottawashooting            &   890   & 470 &   420   & 52.8 \\
putinmissing              &   238   & 126 &   112   & 52.9 \\
sydneysiege$^{\dagger}$   & 1{,}221 & 522 &   699   & 42.8 \\
\midrule
\textbf{Total}            & \textbf{6{,}178} & \textbf{2{,}159} & \textbf{4{,}019} & \textbf{34.9} \\
\bottomrule
\end{tabular}
\end{table}

\subsection{Graph Construction and Node Features}

\paragraph{Graph construction.}
We parse each PHEME thread directory into a directed graph: the source tweet
becomes the root node (index~0), and each reaction tweet adds a directed edge
from its parent (identified via \texttt{in\_reply\_to\_status\_id}) to itself.
Reactions whose parent is absent from the cascade are connected directly to the
root, a common artefact of Twitter data collection. The result is a rooted
directed tree where edges point from parent to child.

\paragraph{Node features.}
Each node carries a $d = 390$-dimensional feature vector: 384-dim text
embeddings from \texttt{all-MiniLM-L6-v2}~\cite{reimers2019sbert},
2 user features ($\log(1{+}\text{followers})$ and verified status), and
4 structural features (PageRank, degree, in-degree, out-degree centrality).
Two positional features, \texttt{is\_root} and BFS \texttt{depth}, are
appended at load time, bringing the full dimension to $d = 392$. We evaluate
three feature configurations: \textit{full} ($d = 392$), \textit{text-only}
($d = 386$), and \textit{structural-only} ($d = 6$).

\subsection{Evaluation Protocol}

We use a two-phase LOEO design to prevent hyperparameter leakage from the
test events (charliehebdo and sydneysiege, marked $\dagger$ in
Table~\ref{tab:dataset}). In the \textit{search phase} these two events are
withheld entirely; 50 guided trials per model-configuration pair are run over
the four training events with macro F1 as the criterion. In the
\textit{evaluation phase}, models are retrained on all four training events
and evaluated on the 3{,}106-graph combined test set. We report macro F1,
precision, recall, and accuracy, together with bootstrap 95\% confidence
intervals and McNemar's test for pairwise comparisons.

\section{Methodology}

\subsection{Baseline Models}

\paragraph{Random Forest (RF).}
Each cascade is reduced to a single vector by mean-pooling its node feature
matrix: $\mathbf{g} = \frac{1}{|\mathcal{V}|}\sum_{v}\mathbf{x}_v$. A Random
Forest classifier (100 trees, balanced class weights) is trained on these
vectors. Any GNN that outperforms the RF gains something specifically from
message passing rather than from richer pooling.

\paragraph{Multi-Layer Perceptron (MLP).}
The MLP also operates on mean-pooled vectors, using fully connected layers with
ReLU activations and dropout, followed by a single output logit. It shares the
same input as the RF but is differentiable and depth-tunable, isolating the
contribution of a learned non-linear classifier before any graph structure is
introduced.

\subsection{Graph Neural Network Models}

Both GNN architectures share three design choices that go beyond their respective
base convolutions: residual skip connections, layer normalisation, and a root
readout that concatenates the mean-pooled graph embedding with the root node's
final representation.

\paragraph{Graph Convolutional Network (GCN).}
We use the spectral graph convolution of Kipf and Welling~\cite{kipf2017semi}.
At each layer $\ell$, the representation of node $v$ is updated as:
%
\begin{equation}
  \tilde{\mathbf{h}}_v^{(\ell)} =
    \sum_{u \in \mathcal{N}(v) \cup \{v\}}
      \frac{1}{\sqrt{\tilde{d}_v}\sqrt{\tilde{d}_u}}\,
      \mathbf{h}_u^{(\ell-1)}\, \mathbf{W}^{(\ell)}
\end{equation}
%
where $\tilde{d}_v = |\mathcal{N}(v)| + 1$ is the degree after adding self-loops.
The pre-activations are passed through a residual block,
$\mathbf{h}_v^{(\ell)} = \mathrm{LayerNorm}(\mathrm{ReLU}(\tilde{\mathbf{h}}_v^{(\ell)} + \mathbf{W}_{\mathrm{skip}}^{(\ell)}\mathbf{h}_v^{(\ell-1)}))$,
where $\mathbf{W}_{\mathrm{skip}}^{(\ell)}$ is a bias-free linear projection.
LayerNorm stabilises training on the shallow cascades common in PHEME.

\paragraph{Graph Attention Network (GAT).}
We use the GATv2 attention mechanism of Brody et al.~\cite{brody2022gatv2}, which
addresses a rank-collapse limitation of the original
GAT~\cite{velickovic2018gat} by computing a fully dynamic attention score:
%
\begin{equation}
  e_{ij} = \mathbf{a}^\top
    \mathrm{LeakyReLU}\!\left(
      \mathbf{W}\, [\mathbf{h}_i^{(\ell-1)} \,\|\, \mathbf{h}_j^{(\ell-1)}]
    \right)
\end{equation}
%
Scores are normalised via softmax over the neighbourhood,
$\alpha_{ij} = \exp(e_{ij})/\sum_{k \in \mathcal{N}(i)}\exp(e_{ik})$.
With $K$ parallel heads, representations are concatenated and passed through
the same residual block as the GCN, with $\mathbf{W}_{\mathrm{skip}}^{(\ell)}$
projecting to $K \times d_\ell$. The number of heads $K$ is tuned during
hyperparameter search.

\paragraph{Shared readout.}
After $L$ message-passing layers, both models produce a graph-level
representation by concatenating the mean-pooled node embeddings with the root
node's final embedding:
%
\begin{equation}
  \hat{\mathbf{g}} =
    \left[
      \frac{1}{|\mathcal{V}|} \sum_{v \in \mathcal{V}} \mathbf{h}_v^{(L)}
      \;\Big\|\;
      \mathbf{h}_0^{(L)}
    \right]
\end{equation}
%
The root node (index~0, the source tweet) carries information about the cascade's
origin that mean pooling would dilute. A single linear layer maps $\hat{\mathbf{g}}$
to the output logit: $\hat{y} = \sigma(\mathbf{w}_c^\top \hat{\mathbf{g}})$.

\paragraph{Edge direction.}
The raw cascade graphs are directed (parent $\to$ child). We experiment with three
configurations: \textit{original} (edges point from source to reply, information
flows bottom-up through the classifier), \textit{inverted} (reversed), and
\textit{bidirectional} (both directions added via \texttt{to\_undirected},
effectively doubling the edge count). Edge direction is treated as a discrete
hyperparameter during search.

\subsection{Training Procedure}

\paragraph{Optimiser and schedule.}
All neural models are trained with the Adam optimiser and a cosine annealing
learning rate schedule with minimum learning rate $\eta_{\min} = \eta_0 / 10$,
where $\eta_0$ is the initial rate. Gradients are clipped to $\ell_2$ norm~1.0
at every step to prevent instability on very small graphs.

\paragraph{Loss function.}
We use binary focal loss~\cite{lin2017focal}:
%
\begin{equation}
  \mathcal{L} = -\frac{1}{N}\sum_{i=1}^{N}
    (1 - p_{t,i})^\gamma \log p_{t,i}
\end{equation}
%
where $p_{t,i} = \sigma(\hat{y}_i)$ for positive examples and
$p_{t,i} = 1 - \sigma(\hat{y}_i)$ for negative examples. The modulating factor
$(1-p_{t,i})^\gamma$ down-weights well-classified examples and concentrates
learning on hard or misclassified ones, which is particularly beneficial given
the class imbalance present in several events. The exponent $\gamma$ is tuned
during hyperparameter search.

\paragraph{Class balancing.}
Before each training run, the training set is balanced by per-event
majority-class undersampling: within each event, the larger class is randomly
subsampled to match the size of the smaller one. This prevents events with
heavily skewed class ratios from dominating the training signal.

\paragraph{Hyperparameter search.}
Hyperparameters are selected with a sequential model-free search: 50 guided
trials are run per model-configuration combination (12 combinations in total:
4 models $\times$ 3 feature configurations), where each trial evaluates the
full LOEO procedure over the four training events and reports the global macro
F1. The search space covers hidden layer dimensions, dropout rate, learning
rate, weight decay, focal loss exponent $\gamma$, number of attention heads
(GAT only), and edge direction (GCN and GAT). The best configuration per
combination is used for all final held-out evaluations.

\section{Related Work}

\paragraph{GNN-based rumour and fake news detection.}
Graph-based approaches to rumour detection represent information cascades as graphs
and apply message-passing to exploit propagation structure. Bian et
al.~\cite{bian2020bigcn} introduced BiGCN, which propagates representations
bidirectionally along the cascade tree (top-down and bottom-up), consistently
outperforming sequential baselines on Twitter15/16. Lu and Li~\cite{lu2020gcan}
proposed GCAN, which jointly models user interaction graphs and source-content
co-attention, also providing post-hoc explanations via attention weights. Monti et
al.~\cite{monti2019fakenews} applied spectral graph convolutions to the propagation
graphs of FakeNewsNet, demonstrating that geometric deep learning transfers across
datasets. Relative to these works, our study makes three methodological departures:
we use a strict two-phase LOEO protocol that keeps two events entirely hidden
during hyperparameter search; we perform a controlled feature ablation that
isolates the contribution of message passing from that of the text embedding; and
we complement classification with a network-science characterisation of cascade
topology.

\paragraph{Explainability and cascade topology.}
GNNExplainer~\cite{ying2019gnnexplainer} provides post-hoc explanations by
optimising a soft mask over the node features and subgraph of each prediction;
we use it to compare the feature-group importance the model assigns to rumour
versus non-rumour graphs. For cascade topology, Goel et
al.~\cite{goel2016structural} defined structural virality as the average
pairwise distance in a diffusion tree, providing a principled measure of
broadcast versus chain-like spread that we adopt as one of our six topological
analyses.

\section{Network Science Analysis}
\label{sec:topology}

We characterise the structural properties of the full 6{,}425-graph dataset
(all events, no size filter) across six analyses. All comparisons use two-sided Mann-Whitney~U tests. The rank-biserial effect
size $r \in [0, 1]$ measures practical magnitude as the non-parametric
correlation between class membership and the metric (0 = no separation,
1 = perfect separation); thresholds: small $r \geq 0.1$, medium $r \geq 0.3$.
Cascade size and depth additionally report Spearman~$\rho$, the rank-order
correlation between metric value and the binary Rumour/Non-Rumour label. The central finding
is that rumour cascades are consistently shallower and more broadcast-oriented
than non-rumour cascades, but all effects are small ($r < 0.15$), and the
structural signal does not generalise across events (confirmed in
Section~\ref{sec:results}).

\paragraph{Cascade size and depth.}
Rumour cascades are significantly smaller (mean 13.89 nodes vs.\ 18.00;
Spearman $\rho = -0.119$, $p < 0.001$) and shallower (mean depth 1.257
vs.\ 1.528; $\rho = -0.102$, $p < 0.001$). Both effects are statistically
robust at $n = 6{,}425$ but small in magnitude.

\paragraph{Tipping point.}
Stratifying by cascade size reveals a tipping-point pattern: rumour prevalence
declines monotonically from 44.9\% among cascades with 1--5 nodes to 11.1\%
in those with more than 200 nodes, consistent with factual content continuing
to accumulate replies while rumour threads stall early.

\paragraph{Structural virality.}
Following Goel et al.~\cite{goel2016structural}, structural virality is defined
as the average pairwise path length (APL) in the cascade tree. Rumour cascades
have significantly lower APL (mean 2.463 vs.\ 2.808, $p < 0.001$, $r = 0.114$),
confirming a structurally more star-like topology (replies cluster around the root)
while non-rumour cascades exhibit longer information chains consistent with viral
diffusion. The effect is small.

\paragraph{Network robustness and echo chambers.}
Robustness is measured by tracking the largest connected component fraction after
removing the top 1\%--30\% of degree-ranked hubs. Both classes show near-identical
curves, differing by at most 0.014 at any removal rate, with the curves crossing at
10\% removal. The branching factor (mean out-degree of internal nodes) is also
indistinguishable between classes ($p = 0.650$, $r = 0.007$).

\paragraph{Novel topological metrics.}
Table~\ref{tab:topology} reports six cascade complexity metrics. Four discriminate
significantly. Root Dominance (fraction of all edges incident on the root node) is
higher for rumours ($r = 0.074$), reinforcing the broadcast/star pattern. Strahler
Number, a recursive measure of tree branching complexity analogous to stream order in
hydrology, is lower for rumours ($r = 0.104$), indicating structurally simpler cascade
trees. Normalised Depth (max depth divided by $\log_2 n$) is lower for rumours
($r = 0.070$), showing that rumour cascades are shallower than expected for their size.
The Gini coefficient of the out-degree distribution (0 = all nodes have
equal out-degree; 1 = one node holds all edges) is lower for rumours
($r = 0.065$), meaning non-rumour cascades contain more high-degree hub nodes.
Leaf Ratio (fraction of nodes with out-degree zero, i.e., no replies) is
marginally significant ($p = 0.010$, $r = 0.039$). Epidemic~$R_0$, the mean
out-degree of nodes that themselves received at least one reply (analogous to
the basic reproduction number in epidemiology), does not reach significance
($p = 0.081$).

\begin{table}[t]
\caption{Novel topological metrics comparing Rumour (R) and Non-Rumour (NR) cascades
($n = 6{,}425$, Mann-Whitney~U). $r$ = rank-biserial effect size.
Significance: *** $p < 0.001$, ** $p < 0.01$, ns $p \geq 0.05$.}
\label{tab:topology}
\centering
\begin{tabular}{lrrrrl}
\toprule
Metric & R mean & NR mean & $p$-value & $r$ & \\
\midrule
Root Dominance   & 0.765 & 0.727 & $<0.001$ & 0.074 & *** \\
Gini out-degree  & 0.683 & 0.718 & $<0.001$ & 0.065 & *** \\
Strahler Number  & 1.893 & 2.024 & $<0.001$ & 0.104 & *** \\
Leaf Ratio       & 0.736 & 0.722 & 0.010    & 0.039 & **  \\
Norm.\ Depth     & 0.726 & 0.822 & $<0.001$ & 0.070 & *** \\
Epidemic $R_0$   & 4.214 & 4.469 & 0.081    & 0.028 & ns  \\
\bottomrule
\end{tabular}
\end{table}

\section{Classification Results}
\label{sec:results}

We evaluate four classifiers under three feature configurations on the fixed held-out
test set (charliehebdo, 1{,}952 graphs; sydneysiege, 1{,}154 graphs; 3{,}106
combined). The random-classifier lower bound (predict-all-positive at $p = 0.5$)
gives a combined macro F1 of 0.375. All neural models are trained with the
best hyperparameters found during the search phase, which excluded the two test
events entirely.

Table~\ref{tab:results} reports per-event and combined macro F1 for the six
relevant model-configuration pairs.

\begin{table}[t]
\caption{Final held-out evaluation. Models are trained on four events and tested on
charliehebdo and sydneysiege (never seen during hyperparameter search).
Combined macro F1 over 3{,}106 graphs; $\Delta$ is the margin above the
random lower bound (0.375). Best result per column in \textbf{bold}.}
\label{tab:results}
\centering
\begin{tabular}{llrrrrr}
\toprule
Model & Config & Charlie F1 & Sydney F1 & Comb.\ F1 & Acc & $\Delta$ \\
\midrule
GAT  & text   & 0.552 & 0.653 & \textbf{0.610} & \textbf{0.770} & +0.235 \\
GAT  & full   & \textbf{0.557} & 0.576 & 0.565 & 0.709 & +0.190 \\
GCN  & text   & 0.504 & 0.604 & 0.551 & 0.686 & +0.176 \\
GCN  & full   & 0.540 & 0.557 & 0.548 & 0.697 & +0.173 \\
MLP  & text   & 0.473 & 0.577 & 0.520 & 0.643 & +0.145 \\
MLP  & full   & 0.460 & \textbf{0.600} & 0.521 & 0.614 & +0.146 \\
\midrule
MLP  & struct & 0.365 & 0.000 & 0.303 & 0.353 & $-$0.072 \\
GCN  & struct & 0.289 & 0.282 & 0.287 & 0.414 & $-$0.088 \\
GAT  & struct & 0.152 & 0.247 & 0.200 & 0.666 & $-$0.175 \\
\bottomrule
\end{tabular}
\end{table}

\paragraph{Structural-only models.}
All three structural-only classifiers fall below the random lower bound (combined
F1 of 0.200--0.303 vs.\ 0.375). This confirms the finding from
Section~\ref{sec:topology}: the topological differences identified between classes
are real but too small to support cross-event generalisation.

\paragraph{Text and full-feature models.}
All six text and full-feature models exceed the lower bound by comfortable margins
(+0.145 to +0.235). GAT with text features achieves the best combined F1 of
0.610. Adding structural and positional features to the text-based models
consistently hurts generalisation: both GAT (0.610 $\to$ 0.565) and GCN
(0.551 $\to$ 0.548) perform worse in the full configuration.
The structural features overfit the training event distributions and do not
transfer to the two test events.

\paragraph{Statistical tests.}
Bootstrap 95\% CIs for the top four models (GAT-text [0.583, 0.636], GAT-full
[0.540, 0.590], GCN-text [0.526, 0.576], GCN-full [0.521, 0.573]) all exclude
the random lower bound, confirming significance above chance.
All pairwise comparisons on the combined 3{,}106-prediction test set are
significant at $p < 0.001$. Message passing beats flat MLP in both feature
settings, confirming that cascade graph structure provides genuine signal
beyond mean-pooled embeddings. Attention (GAT) beats GCN in both settings,
justifying the added mechanism complexity. Adding structural features
significantly hurts both GCN and GAT, indicating that topology encodes
event-specific rather than event-general patterns. The GNN-MLP F1 margin is
modest (+0.031 for text, +0.027 for full), suggesting NLP embeddings carry
the bulk of the discriminative signal.

\section{XAI Analysis}
\label{sec:xai}

To understand what drives the GNN predictions, we apply two complementary
explainability methods on a separate training run using the charliehebdo event
as the held-out test fold (1{,}952 graphs: 436 Rumour, 1{,}516 Non-Rumour).
The models trained for XAI are not the HP-searched models from
Section~\ref{sec:results}; they use default hyperparameters and are trained on
the remaining events. GNNExplainer~\cite{ying2019gnnexplainer} is applied to the
GCN, and attention analysis is applied to the GAT.

\paragraph{GNNExplainer: feature group importance.}
GNNExplainer~\cite{ying2019gnnexplainer} optimises a soft per-node feature
mask to identify which features drove each prediction. We apply it to 40
balanced test graphs, average mask weights within semantic feature groups, and
compare per-graph group importance between classes via Mann-Whitney~U; higher
importance means the model relied more on that group.
The GCN relies significantly more on every feature group when processing Rumour
graphs than Non-Rumour graphs. The largest
differences are in authority signals: follower count ($r = 0.508$, $p < 0.001$)
and PageRank ($r = 0.513$, $p < 0.001$). The NLP embedding group is also more
important for Rumour graphs ($r = 0.464$, $p < 0.001$), as are degree and
in-degree centrality ($r \approx 0.49$, $p < 0.001$). Out-degree is the only
non-significant group ($p = 0.889$), suggesting the number of replies a user
generates is not discriminative. Rumour cascades are thus flagged by a combination
of anomalous authority signals and unusual text, rather than any single feature.

\paragraph{GAT attention by cascade depth.}
The second GAT layer assigns each edge an attention weight $\alpha_{uv}$
(Eq.~2); we aggregate these by the BFS depth of the receiving node (depth~0
= root, depth~1 = direct replies, etc.) and compare across classes.
Depth-2 nodes receive disproportionately high attention in both classes
(mean 0.719--0.776 vs.\ approximately 0.50 at all other depths). Non-Rumour
cascades receive significantly more depth-2 attention than Rumour cascades
(0.776 vs.\ 0.719, $p < 0.001$, $r = 0.233$), suggesting non-rumour threads
generate richer second-tier engagement. The pattern reverses at depth~3
(Rumour $>$ Non-Rumour, $p = 0.001$, $r = 0.106$); beyond depth~3 differences
are not significant.

\paragraph{Attention entropy.}

Per-graph attention entropy measures how diffuse or concentrated attention is
(higher = spread evenly across nodes; lower = focused on a few nodes).
Rumour cascades have significantly lower entropy (mean 2.394 vs.\ 2.615 for
Non-Rumour, $p < 0.001$), indicating that the GAT concentrates on fewer nodes
when classifying rumours. This aligns with the GNNExplainer finding: rumour
cascades appear to contain a small number of structurally or semantically
anomalous nodes (high-authority users, unusual depth-2 reply patterns) that
dominate the model's decision, whereas non-rumour engagement is more evenly
distributed.

\bibliographystyle{splncs04}
\bibliography{references}

\end{document}